\subsection{Quotients of regular local rings}

PROPOSITION 2. --- Let A be a Noetherian local ring, \(\boldsymbol{x} = (x_1, ..., x_r)\) a sequence of elements of \(\mathfrak{m}_{\mathrm{A}}\) and \(\mathfrak{x}\) is the ideal generated by \(\boldsymbol{x}\). The following conditions are equivalent:

\begin{enumerate}
\def\labelenumi{(\roman{enumi})}
\item
  the ring A is regular, and x is part of a coordinate system of A;
\item
  the ring A/x is regular and x is a secant sequence for A ;
\item
  the ring A/x is regular and x is a completely secant sequence for A.
\end{enumerate}

Moreover, when these conditions are satisfied, x is a prime ideal of A.

\begin{enumerate}
\def\labelenumi{(\roman{enumi})}
\setcounter{enumi}{2}
\item
  \(\Rightarrow\) (ii): this follows from the corollary of Prop. 3 in § 3, no. 2.
\item
  \(\Rightarrow\) (i): suppose that \(x\) is a secant sequence for A and that the Noetherian local ring A/\(x\) is regular. Let \((x_{r+1}, ..., x_d)\) be a sequence of elements of A, whose classes modulo \(x\) form a system of coordinates of A/\(x\). Then the sequence \((x_1, ..., x_d)\) generates the ideal \(m_A\) of A, and we have
\end{enumerate}

\[
\dim (\mathrm{A}) = r + \dim (\mathrm{A} / \mathrm{x}) = r + (d - r) = d.
\]

Consequently, A is regular and \((x_{1},\ldots,x_{d})\) is a system of coordinates of A.

\begin{enumerate}
\def\labelenumi{(\roman{enumi})}
\tightlist
\item
  \(\Rightarrow\) (iii): if condition (i) is satisfied, the sequence \(x\) is completely secant (no. 2, th. 1), hence secant by the corollary of prop. 3 of § 3, no. 2. We therefore have
\end{enumerate}

\[
\dim (\mathrm{A} / \mathfrak {x}) = \dim (\mathrm{A}) - r;\tag{3}
\]

moreover, the classes of \(x_{1},\ldots,x_{r}\) modulo \(\mathfrak{m}_{\mathbf{A}}^{2}\) are linearly independent over the field \(\kappa_{\mathbf{A}}\) , and we therefore have

\[
\left[ \mathfrak {m} _ {\mathrm{A}} / (\mathfrak {m} _ {\mathrm{A}} ^ {2} + x): \kappa_ {\mathrm{A}} \right] = \left[ \mathfrak {m} _ {\mathrm{A}} / \mathfrak {m} _ {\mathrm{A}} ^ {2}: \kappa_ {\mathrm{A}} \right] - r.\tag{4}
\]

Formulas (3) and (4) show that A/x is regular.

Every regular Noetherian local ring is an integral domain by Cor. 1 of Th. 1 of No. 2. Consequently, x is prime if A/x is regular.

COROLLARY 1. --- Let A be a Noetherian local ring, and let t be an element of \(\mathfrak{m}_{\mathrm{A}}\). The following conditions are equivalent:

\begin{enumerate}
\def\labelenumi{(\roman{enumi})}
\item
  A is regular, and t does not belong to \(m_{A}^{2}\) ;
\item
  A/tA is regular and \(\dim(\mathrm{A}/t\mathrm{A}) < \dim(\mathrm{A})\) ;
\item
  A/tA is regular, and t is not a zero divisor in A.
\end{enumerate}

COROLLARY 2. --- Let A be a regular Noetherian local ring, and let q be an ideal of A. Then A/q is regular if and only if q is generated by a subset of a system of coordinates of A.

The condition is sufficient by Prop. 2.

Suppose A/q is regular, and let \(x = (x_{1}, \ldots, x_{r})\) be a sequence of elements of q whose classes modulo \(m_{A}^{2}\) form a basis of \((\mathfrak{q} + \mathfrak{m}_{\mathrm{A}}^{2})/\mathfrak{m}_{\mathrm{A}}^{2}\) over the field \(\kappa_{A}\). Let x be the ideal of A generated by x. We therefore have \(x \subset q\) and x is part of a system of coordinates of A, so the Noetherian local ring A/x is regular (prop. 2); moreover, the vector spaces \(\mathfrak{m}_{\mathrm{A}}/(\mathfrak{q} + \mathfrak{m}_{\mathrm{A}}^{2})\) and \(\mathfrak{m}_{\mathrm{A}}/(\mathfrak{x} + \mathfrak{m}_{\mathrm{A}}^{2})\) have the same dimension over \(\kappa_{A}\). Consequently, the regular Noetherian local rings A/q and A/x have the same dimension. Since the ideals q and x are prime and we have \(x \subset q\) , we finally have q = x.

Example. --- Let \(k\) be a field, \(\mathbf{A} = k[[\mathbf{X}_1, ..., \mathbf{X}_n]]\) and \(\mathfrak{q}\) an ideal of \(\mathbf{A}\) , distinct from \(\mathbf{A}\) . For \(\mathbf{A}/\mathfrak{q}\) to be regular, it is necessary and sufficient that one can find an integer \(r \geqslant 0\) and elements \(\mathbf{F}_1, ..., \mathbf{F}_r\) of \(\mathbf{A}\) , generating \(\mathfrak{q}\) , and such that the matrix \(\left(\frac{\partial\mathbf{F}_i}{\partial\mathbf{X}_j}(0, ..., 0)\right)\) be of rank \(r\) ("Jacobian criterion"). We then have \(\dim(\mathbf{A}/\mathfrak{q}) = n - r\) .

Remark. --- Let A be a regular Noetherian local ring and \(q \subset m_{A}\) an ideal of A such that A/q is regular. Let \((x_{1}, \ldots, x_{r})\) be a sequence of elements of q whose classes modulo \(m_{A}^{2}\) generate the vector space \((\mathfrak{q} + \mathfrak{m}_{\mathrm{A}}^{2}) / \mathfrak{m}_{\mathrm{A}}^{2}\) over the field \(\kappa_{A}\). The proof of cor. 2 shows that the ideal q of A is generated by \((x_{1}, \ldots, x_{r})\) .

